\documentclass[aspectratio=169,11pt]{beamer}
\usetheme{Madrid}
\usecolortheme{whale}

% Hỗ trợ tiếng Việt cho pdflatex
\usepackage[T5]{fontenc}
\usepackage[utf8]{inputenc}
\usepackage[vietnamese]{babel} % Optional, có thể bỏ nếu gây lỗi

\usepackage{amsmath,amssymb}
\usepackage{booktabs}
\usepackage{tikz}
\usetikzlibrary{shapes.geometric,arrows.meta,positioning,fit,backgrounds}
\usepackage{xcolor}
\usepackage{hyperref}

% Custom colors
\definecolor{tragblue}{RGB}{0,82,155}
\definecolor{tragorange}{RGB}{230,115,0}
\definecolor{traggreen}{RGB}{0,140,70}
\definecolor{traggray}{RGB}{90,90,90}
\definecolor{lightblue}{RGB}{220,235,250}
\definecolor{lightorange}{RGB}{255,240,220}

\setbeamercolor{frametitle}{bg=tragblue,fg=white}
\setbeamercolor{title}{fg=white}
\setbeamercolor{structure}{fg=tragblue}
\setbeamercolor{block title}{bg=tragblue,fg=white}
\setbeamercolor{block title alerted}{bg=tragorange,fg=white}
\setbeamercolor{block title example}{bg=traggreen,fg=white}

\title[\textbf{T-RAG}]{\textbf{T-RAG: Targeted Retrieval-Augmented Generation}\\
\large Probabilistic Source-Aware RAG for Enterprise Data}
\author{Trình bày: Đăng Bảo \& Đinh Bằng}
\date{Tháng 7, 2026}
\institute{Research Proposal --- Advisor Meeting}

\begin{document}

%------------------------------------------------
\begin{frame}
\titlepage
\end{frame}

%------------------------------------------------
\begin{frame}{Outline}
\tableofcontents
\end{frame}

%------------------------------------------------
\section{Bối cảnh \& Bài toán}

\begin{frame}{Bài toán: RAG trên dữ liệu Doanh nghiệp}

\begin{columns}
\begin{column}{0.5\textwidth}
\begin{block}{Dataset: EnterpriseRAG-Bench}
\begin{itemize}
    \item \textbf{511,962 documents} từ 9 nguồn khác nhau
    \item Slack (285k), Gmail (121k), Jira (41k),\\Confluence (5k), GitHub, Fireflies...
    \item \textbf{500 câu hỏi} benchmark, 10 loại
    \item Dữ liệu nhiễu (Noisy Data): Xung đột ngữ nghĩa, trùng lặp thông tin
\end{itemize}
\end{block}
\end{column}
\begin{column}{0.5\textwidth}
\begin{alertblock}{3 Vấn đề của Standard RAG}
\begin{enumerate}
    \item \textbf{Sự mơ hồ về nguồn (Source Ambiguity)}: Gây phân tán không gian đặc trưng và tăng độ nhiễu.
    \item \textbf{Bất đồng từ vựng (Vocabulary Mismatch)}: Sự sai lệch giữa ngôn ngữ tự nhiên và thuật ngữ kỹ thuật.
    \item \textbf{Truy vấn đa nguồn (Multi-source Queries)}: Đòi hỏi năng lực suy luận chéo nguồn phức tạp.
\end{enumerate}
\end{alertblock}
\end{column}
\end{columns}

\vspace{0.4cm}
\begin{center}
\textit{Kiến trúc Standard RAG bộc lộ hạn chế về độ phức tạp tính toán và khả năng định chuẩn khi đối sánh trên toàn cục không gian dữ liệu.}
\end{center}
\end{frame}

%------------------------------------------------
\section{Tính Mới \& Đóng góp}

\begin{frame}{3 Đóng góp Học thuật}

\begin{block}{\#1 --- Probabilistic Source Router (PSR) + Database Sharding}
\textbf{System-ML Co-design}: Tích hợp phân vùng hệ quản trị Vector DB và mô hình dự báo xác suất định tuyến, tối ưu hóa không gian tìm kiếm từ \textbf{50--90\%}.
\end{block}

\vspace{0.2cm}
\begin{block}{\#2 --- Source-Weighted RRF (SW-RRF)}
\textbf{Cải tiến mô hình toán học}: Đề xuất tích hợp \textbf{Bayesian Prior} trực tiếp vào công thức Reciprocal Rank Fusion nhằm tối ưu hóa tính điểm xếp hạng tài liệu đa nguồn.
\end{block}

\vspace{0.2cm}
\begin{block}{\#3 --- Cross-Source Entity Propagation (CSEP)}
\textbf{Phương pháp suy luận đa bước}: Khắc phục giới hạn truy vấn chéo bằng kỹ thuật trích xuất thực thể mỏ neo (Anchor Entity) để định hướng chuỗi suy luận thứ cấp.
\end{block}
\end{frame}

%------------------------------------------------
\section{Kiến trúc Pipeline Đề xuất}

\begin{frame}{Tổng quan Pipeline T-RAG (4 Giai đoạn)}

\resizebox{\textwidth}{!}{
\begin{tikzpicture}[
  node distance=0.6cm and 0.4cm,
  box/.style={rectangle, rounded corners=4pt, minimum width=2.5cm, minimum height=0.9cm, text centered, font=\small\bfseries, text width=2.5cm},
  arrow/.style={-{Stealth[length=4pt]}, thick},
  bluebox/.style={box, fill=tragblue!20, draw=tragblue, text=tragblue},
  orangebox/.style={box, fill=tragorange!20, draw=tragorange, text=tragorange},
  greenbox/.style={box, fill=traggreen!20, draw=traggreen, text=traggreen},
  graybox/.style={box, fill=gray!15, draw=gray, text=black},
]

\node[graybox] (q) {User Query $Q$};
\node[bluebox, right=1.2cm of q] (psr) {Giai đoạn 1\\PSR Router};
\node[bluebox, right=1.2cm of psr] (shard) {DB Sharding\\Sub-space};
\node[orangebox, right=1.2cm of shard] (hybrid) {Giai đoạn 2\\SW-RRF\\Hybrid Search};
\node[orangebox, right=1.2cm of hybrid] (csep) {Giai đoạn 3\\CSEP\\Multi-hop};
\node[greenbox, below=1.0cm of csep] (rerank) {Giai đoạn 4\\Cross-Encoder\\Reranker};
\node[greenbox, left=1.2cm of rerank] (gen) {LLM\\Generation};
\node[graybox, left=1.2cm of gen] (ans) {Answer};

\draw[arrow] (q) -- (psr);
\draw[arrow] (psr) -- node[above,font=\scriptsize]{$P(s_i|Q)$} (shard);
\draw[arrow] (shard) -- node[above,font=\scriptsize]{$S_{active}$} (hybrid);
\draw[arrow] (hybrid) -- node[above,font=\scriptsize]{candidates} (csep);
\draw[arrow] (csep) -- node[right,font=\scriptsize,text width=1.8cm,align=center]{multi-hop\\docs} (rerank);
\draw[arrow] (rerank) -- node[above,font=\scriptsize]{top-K} (gen);
\draw[arrow] (gen) -- (ans);

\node[font=\tiny\itshape, below=0.15cm of psr, text=traggray] {Semantic Routing};
\node[font=\tiny\itshape, below=0.15cm of hybrid, text=traggray] {Dense + BM25};
\node[font=\tiny\itshape, below=0.15cm of rerank, text=traggray] {Threshold $\theta$};

\end{tikzpicture}
}

\vspace{0.2cm}
\begin{columns}
\begin{column}{0.33\textwidth}
\centering\small
\textcolor{tragblue}{\textbf{Xanh = Pre-Retrieval}}\\Phân loại \& Định tuyến
\end{column}
\begin{column}{0.33\textwidth}
\centering\small
\textcolor{tragorange}{\textbf{Cam = Retrieval}}\\Lai ghép \& Suy luận đa bước
\end{column}
\begin{column}{0.33\textwidth}
\centering\small
\textcolor{traggreen}{\textbf{Xanh lá = Post-Retrieval}}\\Kiểm soát ảo giác
\end{column}
\end{columns}
\end{frame}

%------------------------------------------------
\begin{frame}{Giai đoạn 1: Probabilistic Source Router (PSR)}
\small
\begin{columns}[T]
\begin{column}{0.54\textwidth}
\textbf{Toán học định tuyến:} 
Dự báo xác suất có điều kiện của nguồn:
\vspace{-0.2cm}
{\footnotesize $$\boxed{P(s_i | Q) = \text{Softmax}(W \cdot \text{Encoder}(Q) + b)_i}$$}

\textbf{Thuật toán Phân vùng (Sharding):}
\begin{itemize}
    \item Chỉ kích hoạt tìm kiếm trên tập nguồn con:
\vspace{-0.2cm}
    {\footnotesize $$S_{active} = \{ s_i \in S \mid P(s_i | Q) \ge \tau \}$$}
\end{itemize}

\vspace{0.1cm}
\textbf{Kết quả:} 
Độ phức tạp giảm từ $O(|D_{total}|)$ xuống $O(\sum_{i \in S_{active}} |D_{s_i}|)$. Giảm I/O đáng kể.
\end{column}
\begin{column}{0.44\textwidth}
\begin{exampleblock}{Ví dụ minh họa định tuyến}
Truy vấn: \textit{"What did the team say about the GPU budget meeting?"}

\vspace{0.1cm}
Kết quả phân phối xác suất từ PSR:
\begin{itemize}
    \item $P(\text{Fireflies}|Q) = 0.81 > \tau \rightarrow$ Kích hoạt
    \item $P(\text{Slack}|Q) = 0.72 > \tau \rightarrow$ Kích hoạt
    \item $P(\text{Jira}|Q) = 0.04 \le \tau \rightarrow$ Loại bỏ
    \item $P(\text{GitHub}|Q) = 0.02 \le \tau \rightarrow$ Loại bỏ
\end{itemize}
\end{exampleblock}
\end{column}
\end{columns}
\end{frame}

%------------------------------------------------
\begin{frame}{Giai đoạn 2: Source-Weighted RRF (SW-RRF)}
\small
\begin{columns}[T]
\begin{column}{0.5\textwidth}
\textbf{RRF truyền thống (Cormack, 2009):}
Đánh đồng mọi tài liệu, làm mất đi sự "ưu tiên" cho nguồn chính.

\vspace{0.2cm}
\textbf{Đề xuất tính mới (SW-RRF):}
Đưa xác suất $P(s_d|Q)$ làm \textbf{Bayesian Prior} vào công thức xếp hạng.

\vspace{-0.2cm}
{\footnotesize $$\boxed{Score(d) = P(s_d|Q)^\gamma \cdot \left( \frac{\alpha}{k + r_{dense}} + \frac{1-\alpha}{k + r_{sparse}} \right)}$$}

\begin{itemize}
    \item $\gamma$: Source Bias Factor
    \item $\alpha$: Trọng số Dense/Sparse
\end{itemize}
\end{column}
\begin{column}{0.48\textwidth}
\begin{block}{Cơ sở lý thuyết của SW-RRF}
Tối ưu hóa thứ hạng tài liệu mục tiêu thông qua trọng số phân phối xác suất ($P(s_d|Q)$), triệt tiêu nhiễu từ các văn bản có tính tương đồng bề mặt nhưng sai lệch ngữ cảnh nguồn.
\end{block}

\vspace{0.1cm}
\begin{exampleblock}{Khả năng bổ trợ của Hybrid Search}
\begin{itemize}
    \item \textbf{Dense Search}: Trích xuất và đối sánh biểu diễn ngữ nghĩa trừu tượng.
    \item \textbf{Sparse (BM25)}: Ràng buộc tính chính xác từ vựng (Exact match).
\end{itemize}
\end{exampleblock}
\end{column}
\end{columns}
\end{frame}

%------------------------------------------------
\begin{frame}{Giai đoạn 3 \& 4: Multi-hop \& Reranking}
\small
\begin{columns}[T]
\begin{column}{0.52\textwidth}
\begin{block}{Giai đoạn 3: CSEP (Suy luận đa bước)}
Mô phỏng cơ chế đồ thị hai phía (Bipartite Graph Random Walk).
\begin{enumerate}
    \item \textbf{Hop 1:} Khởi tạo truy xuất trên nguồn tối ưu nhất $\rightarrow$ Định vị tài liệu mỏ neo ($D_{anchor}$).
    \item \textbf{Entity Extraction:} Phân tách và nhận diện tập thực thể $E$ từ $D_{anchor}$.
    \item \textbf{Hop 2:} Mở rộng không gian truy vấn $Q^{(2)} = Q \oplus E$ để đối sánh chéo nguồn.
\end{enumerate}
\end{block}
\end{column}
\begin{column}{0.46\textwidth}
\begin{block}{Giai đoạn 4: Kiểm soát ngưỡng tin cậy}
Sử dụng \textbf{Cross-Encoder} tối ưu hóa hàm đánh giá $CE(Q, d_{final})$.

\vspace{0.1cm}
Áp dụng \textbf{Hard Thresholding}: Lọc nhiễu triệt để nếu $CE(Q, d) < \theta$.

\vspace{0.1cm}
Trong điều kiện $D_{final} = \emptyset$, hệ thống ngắt luồng LLM và trả về "Unanswerable", hỗ trợ giảm thiểu hiện tượng ảo giác (Hallucination).
\end{block}
\end{column}
\end{columns}
\end{frame}

%------------------------------------------------
\section{Kế hoạch Thực thi Kỹ thuật}

\begin{frame}{Kế hoạch Thực thi \& Đánh giá}

\begin{block}{Các bước Triển khai}
\begin{enumerate}
    \item \textbf{Tiền xử lý Dữ liệu}: Áp dụng mô hình Sharding chia 511k documents thành các bảng phân mảnh độc lập theo \texttt{source\_type}.
    \item \textbf{Huấn luyện PSR Router}: Tinh chỉnh (Fine-tune) mô hình Encoder gọn nhẹ thành bộ phân loại xác suất (Soft Classifier).
    \item \textbf{Triển khai SW-RRF}: Xây dựng hàm lai ghép truy xuất (Retriever) tích hợp công thức toán học hệ số.
\end{enumerate}
\end{block}

\vspace{0.3cm}
\begin{block}{Đánh giá (So với Standard RAG)}
\begin{itemize}
    \item \textbf{Chất lượng:} Recall@K, NDCG@10 (sử dụng tập ground truth)
    \item \textbf{Tốc độ/Chi phí:} Latency (ms), Search Space Size (\% documents được load)
\end{itemize}
\end{block}
\end{frame}

%------------------------------------------------
\begin{frame}[allowframebreaks]{Tài liệu Tham khảo (Verified References)}
\small
\begin{enumerate}
    \item Qiao et al. (2024). \textit{Route Before Retrieve: A New Paradigm for Retrieval-Augmented Generation}.
    \item Wang et al. (2024). \textit{RAGRouter: Learning to Route in Retrieval-Augmented Generation}. NeurIPS 2024.
    \item Cormack, G. V., Clarke, C. L., \& Büttcher, S. (2009). \textit{Reciprocal rank fusion outperforms Condorcet and individual rank learning methods}. ACM SIGIR.
    \item Tang et al. (2024). \textit{MultiHop-RAG: Benchmarking Retrieval-Augmented Generation for Multi-Hop Queries}. arXiv:2401.15391.
    \item Zhang et al. (2025). \textit{HopRAG: Multi-Hop Reasoning for Logic-Aware Retrieval-Augmented Generation}. arXiv:2502.12442.
    \item Nogueira, R., \& Cho, K. (2019). \textit{Passage Re-ranking with BERT}. arXiv:1901.04085.
\end{enumerate}
\end{frame}

\end{document}